Un satèl·lit descriu una òrbita el·líptica al voltant de la Terra. L'energia mecànica d'aquest satèl·lit és $E = -\frac{GM_Tm}{2a}$ on $a$ és la longitud del semieix major, $M_T$ és la massa de la Terra i $m$ la massa del satèl·lit. Siguin $r_p$ el radi del perigeu i $r_a$ el radi de l'apogeu i $v_p$ la velocitat al perigeu.

\begin{apartat}{1,25}
Demostreu que l'energia mecànica del satèl·lit es pot expressar com:
\[ E = -\frac{GM_TL_{S,T}}{r_p(r_p + r_a)v_p}, \]
on $L_{S,T}$ és el mòdul del moment angular del satèl·lit respecte el centre de la Terra.
\end{apartat}

\begin{apartat}{1,25}
La massa del satèl·lit és 750~kg, té una velocitat de 5,70~km/s al perigeu, el perigeu està situat a 18\,400~km del centre de la Terra i l'apogeu es troba a una distància de 54\,900~km del centre de la Terra. Trobeu la velocitat a l'apogeu i l'energia mecànica del satèl·lit.
\end{apartat}

\dades{%
  $G = 6,67\cdot 10^{-11}\ \mathrm{N\,m^2\,kg^{-2}}$\\
  Massa del la Terra $M_\mathrm{T} = 5,97\cdot 10^{24}\ \mathrm{kg}$}
